\documentclass[9pt]{article}
\usepackage[utf8]{inputenc}
\usepackage{amsmath}
\usepackage{geometry}
\geometry{a4paper, margin=0.4in}

\begin{document}

\begin{center}
    \textbf{\large GUÍA DE RESPUESTAS Y SOLUCIONARIO (PRIMER PARCIAL)} \\
    \textbf{FÍSICA BÁSICA (FIS-0001) -- TEMAS A Y B} \\
    \textbf{INSTITUTO TECNOLÓGICO SUPERIOR COMUNITARIO (ITSC)} \\
    \textit{Escala de Puntuación: Parte I (Teoría) = 4 Puntos (0.5 pto c/u) | Parte II (Ejercicios) = 6 Puntos (1.5 ptos c/u) | Total = 10 Puntos}
\end{center}

\vspace{0.1cm}
\hrule
\vspace{0.2cm}

\noindent
\begin{minipage}[t]{0.48\textwidth}
    \centering \textbf{\large TEMA A (SOLUCIONES)} \\
    \vspace{0.2cm}
    \flushleft
    \textbf{PARTE I: Respuestas Teóricas (4 Puntos)}
    \begin{enumerate}
        \item \textbf{b} (Escribir de forma compacta potencias de base 10).
        \item \textbf{b} (Ceros a la derecha tras el decimal cuentan: 7, 0, 3, 0 $\rightarrow$ 4 cifras).
        \item \textbf{b} (Suma: $13.75 \rightarrow 13.8\text{ g}$. Redondeado a un decimal debido a $12.3$).
        \item \textbf{c} (La fuerza es vectorial; temperatura y masa son escalares).
        \item \textbf{b} (Un vector requiere módulo, dirección y sentido).
        \item \textbf{b} (Teorema de Pitágoras por ser perpendiculares: $R = \sqrt{A_x^2 + A_y^2}$).
        \item \textbf{b} (El desplazamiento es el vector en línea recta inicio-fin).
        \item \textbf{c} (M.R.U.: velocidad constante, aceleración nula).
    \end{enumerate}

    \vspace{0.2cm}
    \textbf{PARTE II: Justificación de Ejercicios (6 Puntos)}
    \begin{enumerate}
        \item \textbf{Notación Científica:} \\
        Operación: $(4.00 \times 10^4\text{ m/s}) \times (2.00 \times 10^{-2}\text{ s})$ \\
        Fórmula: $(A \times 10^a) \times (B \times 10^b) = (A \cdot B) \times 10^{a+b}$ \\
        Cálculo: $(4.00 \cdot 2.00) \times 10^{4 - 2} = 8.00 \times 10^2$ \\
        \textbf{Resultado: } $8.00 \times 10^2\text{ m}$ (o $800.00\text{ m}$)
        
        \item \textbf{Suma de Cifras Significativas:} \\
        Operación: $12.35\text{ kg} + 4.2\text{ kg} + 0.156\text{ kg}$ \\
        Cálculo aritmético: $16.706$ \\
        Regla: El resultado se limita por el sumando con menor cantidad de posiciones decimales (en este caso, $4.2$ tiene una sola cifra decimal). Redondeando $16.706$ a un decimal da $16.7$. \\
        \textbf{Resultado: } $16.7\text{ kg}$
        
        \item \textbf{Componentes de un Vector:} \\
        Datos: $F = 20.00\text{ N}, \, \theta = 30.00^\circ$ \\
        Componente horizontal: $F_x = F \cos\theta$ \\
        $F_x = 20.00 \times \cos 30.00^\circ = 20.00 \times 0.87 = 17.40\text{ N}$ \\
        Componente vertical: $F_y = F \operatorname{sen}\theta$ \\
        $F_y = 20.00 \times \operatorname{sen} 30.00^\circ = 20.00 \times 0.50 = 10.00\text{ N}$ \\
        \textbf{Resultado: } $F_x = 17.40\text{ N} \quad | \quad F_y = 10.00\text{ N}$
        
        \item \textbf{Movimiento Rectilíneo Uniforme (M.R.U.):} \\
        Datos: $v = 25.00\text{ m/s}, \, t = 2.00\text{ min} = 120.00\text{ s}$ \\
        Fórmula: $d = v \cdot t$ \\
        Sustitución: $d = 25.00\text{ m/s} \times 120.00\text{ s} = 3000.00\text{ m}$ \\
        \textbf{Resultado: } $3000.00\text{ m}$ (o $3.00 \times 10^3\text{ m}$)
    \end{enumerate}
\end{minipage}
\hfill
\vrule
\hfill
\begin{minipage}[t]{0.48\textwidth}
    \centering \textbf{\large TEMA B (SOLUCIONES)} \\
    \vspace{0.2cm}
    \flushleft
    \textbf{PARTE I: Respuestas Teóricas (4 Puntos)}
    \begin{enumerate}
        \item \textbf{a} (Escribir cantidades grandes o pequeñas de forma compacta).
        \item \textbf{b} (Ceros a la derecha tras el decimal cuentan: 9, 0, 4, 0 $\rightarrow$ 4 cifras).
        \item \textbf{b} (Resta: $21.42 \rightarrow 21.4\text{ m}$. Redondeado a un decimal debido a $24.6$).
        \item \textbf{c} (La energía es escalar; el peso y el desplazamiento son vectoriales).
        \item \textbf{b} (La dirección de un vector la da el ángulo con el eje horizontal positivo).
        \item \textbf{b} (Teorema de Pitágoras por ser perpendiculares: $R = \sqrt{v_x^2 + v_y^2}$).
        \item \textbf{a} (La distancia mide la longitud del camino real recorrido).
        \item \textbf{b} (La pendiente en $x$ vs $t$ representa directamente la velocidad).
    \end{enumerate}

    \vspace{0.2cm}
    \textbf{PARTE II: Justificación de Ejercicios (6 Puntos)}
    \begin{enumerate}
        \item \textbf{Notación Científica:} \\
        Operación: $(2.00 \times 10^6\text{ m/s}) \times (3.00 \times 10^{-3}\text{ s})$ \\
        Fórmula: $(A \times 10^a) \times (B \times 10^b) = (A \cdot B) \times 10^{a+b}$ \\
        Cálculo: $(2.00 \cdot 3.00) \times 10^{6 - 3} = 6.00 \times 10^3$ \\
        \textbf{Resultado: } $6.00 \times 10^3\text{ m}$ (o $6000.00\text{ m}$)
        
        \item \textbf{Suma de Cifras Significativas:} \\
        Operación: $14.28\text{ s} + 3.3\text{ s} + 0.245\text{ s}$ \\
        Cálculo aritmético: $17.825$ \\
        Regla: El resultado se limita por el sumando con menor cantidad de posiciones decimales (en este caso, $3.3$ tiene una sola cifra decimal). Redondeando $17.825$ a un decimal da $17.8$. \\
        \textbf{Resultado: } $17.8\text{ s}$
        
        \item \textbf{Componentes de un Vector:} \\
        Datos: $v = 10.00\text{ m/s}, \, \theta = 45.00^\circ$ \\
        Componente horizontal: $v_x = v \cos\theta$ \\
        $v_x = 10.00 \times \cos 45.00^\circ = 10.00 \times 0.71 = 7.10\text{ m/s}$ \\
        Componente vertical: $v_y = v \operatorname{sen}\theta$ \\
        $v_y = 10.00 \times \operatorname{sen} 45.00^\circ = 10.00 \times 0.71 = 7.10\text{ m/s}$ \\
        \textbf{Resultado: } $v_x = 7.10\text{ m/s} \quad | \quad v_y = 7.10\text{ m/s}$
        
        \item \textbf{Movimiento Rectilíneo Uniforme (M.R.U.):} \\
        Datos: $v = 15.00\text{ m/s}, \, t = 4.00\text{ min} = 240.00\text{ s}$ \\
        Fórmula: $d = v \cdot t$ \\
        Sustitución: $d = 15.00\text{ m/s} \times 240.00\text{ s} = 3600.00\text{ m}$ \\
        \textbf{Resultado: } $3600.00\text{ m}$ (o $3.60 \times 10^3\text{ m}$)
    \end{enumerate}
\end{minipage}

\end{document}
